\documentclass[11pt,a4paper]{article}
\usepackage[utf8]{inputenc}
\usepackage[vietnamese]{babel}
\usepackage{amsmath,amssymb}
\usepackage{geometry}
\geometry{a4paper, total={180mm,260mm}, left=15mm, right=15mm, top=18mm, bottom=18mm}
\usepackage{tikz}
\usetikzlibrary{shapes.geometric, arrows.meta, positioning, calc}
\usepackage{tcolorbox}
\tcbuselibrary{skins,breakable}
\usepackage{xcolor}
\usepackage{multicol}
\usepackage{array}
\usepackage{fontawesome}
\usepackage{fancyhdr}

\pagestyle{fancy}
\fancyhf{}
\renewcommand{\headrulewidth}{0.5pt}
\renewcommand{\footrulewidth}{0.5pt}
\fancyhead[L]{\small\bfseries\color{teal}\faBook\ SỔ TAY KIẾN THỨC TOÁN 3}
\fancyhead[R]{\small\bfseries\color{teal}Lớp Toán Thầy Lĩnh -- loptoanthaylinh.vn}
\fancyfoot[C]{\small\color{gray}-- Sổ tay kiến thức: Bài 36. Nhân số có ba chữ số với số có một chữ số --}

\definecolor{primary}{HTML}{005F73}
\definecolor{secondary}{HTML}{0A9396}
\definecolor{accent}{HTML}{EE9B00}
\definecolor{danger}{HTML}{AE2012}
\definecolor{boxbg}{HTML}{F4F6F8}

\begin{document}

% --- BANNER TIÊU ĐỀ ---
\begin{tcolorbox}[
  enhanced,
  colback=white,
  colframe=primary,
  arc=3mm,
  boxrule=1.2pt,
  left=4mm, right=4mm, top=2.5mm, bottom=2.5mm,
  drop shadow=black!10!white
]
  \centering
  {\Large\bfseries\color{primary}\faBookmark\ BÀI 36: NHÂN SỐ CÓ BA CHỮ SỐ VỚI SỐ CÓ MỘT CHỮ SỐ}\\[1.5mm]
  {\small\itshape\color{black!75}Kĩ thuật đặt tính, thứ tự nhân từ phải sang trái, nhân không nhớ và có nhớ 1 lần trên 1 trang A4}
\end{tcolorbox}

\vspace{-2mm}


\noindent
\begin{minipage}[t]{0.49\textwidth}\vspace{0pt}
\begin{tcolorbox}[
  enhanced,
  colback=blue!4!white,
  colframe=blue!75!black,
  arc=2mm, boxrule=0.8pt,
  left=2mm, right=2mm, top=1.8mm, bottom=1.8mm,
  title={\footnotesize\bfseries\color{white}\faTimes\ 1. Quy trình Đặt tính rồi tính}
]
{\footnotesize
\begin{itemize}\setlength{\itemsep}{1.5pt}\setlength{\leftskip}{-2.5mm}
  \item \textbf{Đặt tính:} Viết thừa số có 1 chữ số dưới thừa số có 3 chữ số, thẳng cột hàng đơn vị. Viết dấu $\times$, kẻ vạch ngang.
  \item \textbf{Thực hiện phép nhân:} Nhân lần lượt thừa số thứ hai với từng chữ số của thừa số thứ nhất theo thứ tự từ \textbf{phải sang trái}:
  \begin{enumerate}\setlength{\itemsep}{0.5pt}\setlength{\leftskip}{-2.5mm}
    \item Hàng đơn vị trước.
    \item Hàng chục tiếp theo.
    \item Hàng trăm sau cùng.
  \end{enumerate}
\end{itemize}
}
\end{tcolorbox}
\end{minipage}
\hfill
\begin{minipage}[t]{0.49\textwidth}\vspace{0pt}
\begin{tcolorbox}[
  enhanced,
  colback=teal!4!white,
  colframe=teal!75!black,
  arc=2mm, boxrule=0.8pt,
  left=2mm, right=2mm, top=1.8mm, bottom=1.8mm,
  title={\footnotesize\bfseries\color{white}\faKey\ 2. Kĩ thuật Nhớ chính xác}
]
{\footnotesize
\begin{itemize}\setlength{\itemsep}{1.5pt}\setlength{\leftskip}{-2.5mm}
  \item \textbf{Phép nhân không nhớ:} Tích các hàng đều $< 10$, viết trực tiếp kết quả.
  \item \textbf{Phép nhân có nhớ $1$ lần:}
  \begin{itemize}\setlength{\itemsep}{0.5pt}\setlength{\leftskip}{-2.5mm}
    \item \textit{Nhớ sang hàng chục:} Lấy tích hàng đơn vị, viết chữ số đơn vị, nhớ số chục sang tích hàng chục.
    \item \textit{Nhớ sang hàng trăm:} Lấy tích hàng chục cộng số nhớ, nhớ tiếp sang hàng trăm nếu $\ge 10$.
  \end{itemize}
  \item Luôn \textbf{nhân trước rồi mới cộng số nhớ}!
\end{itemize}
}
\end{tcolorbox}
\end{minipage}

\vspace{-2.5mm}
% --- KHỐI 3: SƠ ĐỒ TƯ DUY TIKZ ---
\noindent
\begin{tcolorbox}[
  enhanced,
  colback=orange!4!white,
  colframe=orange!80!black,
  arc=2mm, boxrule=0.8pt,
  left=2mm, right=2mm, top=1.5mm, bottom=1.5mm,
  title={\footnotesize\bfseries\color{white}\faSitemap\ SƠ ĐỒ TƯ DUY: TIẾN TRÌNH NHÂN 3 CHỮ SỐ VỚI 1 CHỮ SỐ}
]
\centering
\begin{tikzpicture}[
  box/.style={draw, rounded corners=2pt, align=center, font=\scriptsize\bfseries, inner sep=2.5pt, minimum height=0.75cm},
  arr/.style={-{Stealth[length=2mm]}, thick}
]
  \node[box, fill=purple!15, draw=purple!80!black, text width=3.2cm] (n1) at (0, 0) {BƯỚC 1: HÀNG ĐƠN VỊ\\$c \times k$\\Viết đơn vị, nhớ chục};
  \node[box, fill=teal!15, draw=teal!80!black, text width=3.4cm] (n2) at (4.5, 0) {BƯỚC 2: HÀNG CHỤC\\$b \times k + \text{nhớ}$\\Viết chữ số chục, nhớ trăm};
  \node[box, fill=blue!12, draw=blue!80!black, text width=3.4cm] (n3) at (9.2, 0) {BƯỚC 3: HÀNG TRĂM\\$a \times k + \text{nhớ}$\\Viết toàn bộ kết quả};

  \draw[arr, purple!80!black] (n1) -- (n2);
  \draw[arr, teal!80!black] (n2) -- (n3);
\end{tikzpicture}
\end{tcolorbox}

\vspace{-2.5mm}
% --- KHỐI 4: VÍ DỤ & MẸO TRÁNH BẪY ---
\noindent
\begin{minipage}[t]{0.49\textwidth}\vspace{0pt}
\begin{tcolorbox}[
  enhanced,
  colback=green!4!white,
  colframe=green!60!black,
  arc=2mm, boxrule=0.8pt,
  left=2mm, right=2mm, top=1.8mm, bottom=1.8mm,
  title={\footnotesize\bfseries\color{white}\faCheckCircle\ Ví dụ thực hành chuẩn mực}
]
{\footnotesize
\textbf{1. Phép nhân không nhớ:}\\
$\begin{array}{r@{}l}
  123 & \\
  \times\quad 3 & \\
  \hline
  \mathbf{369} &
\end{array}$
\qquad
\begin{minipage}[c]{0.55\linewidth}
$3 \times 3 = 9$, viết $9$.\\
$3 \times 2 = 6$, viết $6$.\\
$3 \times 1 = 3$, viết $3$.
\end{minipage}

\textbf{2. Phép nhân có nhớ sang hàng chục:}\\
$\begin{array}{r@{}l}
  216 & \\
  \times\quad 4 & \\
  \hline
  \mathbf{864} &
\end{array}$
\qquad
\begin{minipage}[c]{0.55\linewidth}
$4 \times 6 = 24$, viết $4$ nhớ $2$.\\
$4 \times 1 = 4$, thêm $2$ là $6$, viết $6$.\\
$4 \times 2 = 8$, viết $8$.
\end{minipage}
}
\end{tcolorbox}
\end{minipage}
\hfill
\begin{minipage}[t]{0.49\textwidth}\vspace{0pt}
\begin{tcolorbox}[
  enhanced,
  colback=red!4!white,
  colframe=red!70!black,
  arc=2mm, boxrule=0.8pt,
  left=2mm, right=2mm, top=1.8mm, bottom=1.8mm,
  title={\footnotesize\bfseries\color{white}\faExclamationTriangle\ Mẹo tránh bẫy trắc nghiệm}
]
{\footnotesize
\begin{itemize}\setlength{\itemsep}{1.5pt}\setlength{\leftskip}{-2.5mm}
  \item \textbf{Bẫy quên cộng số nhớ:}\\
  Học sinh thường nhớ trong đầu nhưng khi nhân xong hàng chục lại quên cộng thêm số nhớ. Hãy ghi nhỏ số nhớ bên cạnh để không bị sót!
  \item \textbf{Bẫy cộng nhớ trước khi nhân:}\\
  Tuyệt đối không được lấy $(1 + 2) \times 4$! Quy tắc bất di bất dịch là: \textbf{Nhân trước, cộng nhớ sau}.
\end{itemize}
}
\end{tcolorbox}
\end{minipage}


\end{document}
